\subsection{Simuladores considerados}

Las alternativas evaluadas corresponden a herramientas pertinentes para simulación
cuántica de estado completo, aunque con énfasis distintos en cuanto a rendimiento,
portabilidad, facilidad de uso o experimentación.

\begin{itemize}
    \item \textbf{QuEST (Quantum Exact Simulation Toolkit):} simulador de alto
    rendimiento orientado a vectores de estado y matrices de densidad, con soporte
    para ejecución multihilo, distribuida y acelerada por GPU. Su diseño privilegia
    escalabilidad y desempeño en arquitecturas heterogéneas. \cite{Jones2019, quest_repo}

    \item \textbf{Intel-QS (Intel Quantum Simulator):} simulador de circuitos cuánticos
    basado en representación completa del estado, concebido para aprovechar
    arquitecturas multinúcleo y multinodo. Su enfoque está orientado a entornos HPC y
    a la eficiencia en ejecución paralela. \cite{intelqs_doc}

    \item \textbf{qsim:} biblioteca de simulación de estado completo integrada al
    ecosistema de Cirq, optimizada para calcular el estado final de circuitos cuánticos.
    Su orientación principal está asociada a la ejecución eficiente de circuitos y a su uso
    como backend de simulación. \cite{qsim_repo, qsim_cirq}

    \item \textbf{Quantum++:} biblioteca general en C++ para simulación cuántica,
    distribuida como librería \emph{header-only} y diseñada con énfasis en portabilidad,
    facilidad de uso y soporte multihilo. Su arquitectura la hace adecuada para
    prototipado y experimentación ligera. \cite{qpp_paper}

    \item \textbf{TMFQSfullstate:} prototipo de simulador desarrollado en C++ con foco
    en estudiar estrategias de representación y consumo de memoria sobre simulación
    de estado completo. Su diseño favorece la intervención directa sobre la
    representación del estado y la exploración de variantes relacionadas con gestión de
    memoria y paralelismo. \cite{diaz2024software, diaz_mdpi_mmss, tmfqs_repo}
\end{itemize}

En conjunto, las alternativas cubren dos perfiles. Por un lado, QuEST, Intel-QS y qsim
representan herramientas consolidadas con fuerte orientación a rendimiento y
ejecución eficiente. Por otro, Quantum++ y TMFQSfullstate ofrecen entornos más
favorables para intervención directa sobre el código y experimentación sobre la
representación del estado. Esta distinción es relevante porque la selección debía
responder al problema específico de integración y evaluación planteado en esta tesis.
